\chapter{Methodology: Spatiotemporal Graph Modelling}
\label{ch:methodology}

%==============================================================================
\section{Architecture Overview}
\label{sec:architecture}

The pipeline is deliberately staged rather than end-to-end, and the staging is itself a
hypothesis about the sample-size regime (Section~\ref{sec:setting}) rather than an
inherited convention.

\begin{align*}
  \underbrace{\text{BOLD}}_{\text{4-D volume}}
  \;&\longrightarrow\;
  \underbrace{\mathbf{C}^{(t)} \in \mathbb{R}^{200\times200}}_{\text{FC matrix per visit}}
  \;\longrightarrow\;
  \underbrace{G^{(t)} = (X^{(t)}, A^{(t)})}_{\text{graph per visit}} \\[0.4em]
  \;&\longrightarrow\;
  \underbrace{z^{(t)} \in \mathbb{R}^{d}}_{\text{visit embedding}}
  \;\longrightarrow\;
  \underbrace{h^{(T)}}_{\text{trajectory state}}
  \;\longrightarrow\;
  \underbrace{\hat{y}}_{\text{conversion logit}}
\end{align*}

Stage~1 learns the per-visit map $G^{(t)} \mapsto z^{(t)}$ without labels, on a
reconstruction objective. Stage~2 freezes that map and learns only the trajectory
classifier $\{z^{(t)}\} \mapsto \hat{y}$ from the small labelled set. The intended benefit
is that the expensive representation learning consumes unlabelled data while the scarce
labels are spent only on a small number of parameters. Whether that benefit materialises is
exactly what the ablation of Section~\ref{sec:encoder-arms} tests, and the answer is not
the assumed one.

%==============================================================================
\section{Graph Construction}
\label{sec:graph-construction}

For each subject-visit, the Fisher $z$-transformed Pearson correlation matrix
$\mathbf{C} \in \mathbb{R}^{200\times200}$ is converted to a graph as follows.

\paragraph{Node features.} $X = \mathbf{C}$. Node $i$'s feature vector is row $i$ of the
connectivity matrix --- that ROI's correlation profile against all 200 regions. The node
features are therefore the connectivity data in full, at full precision.

\paragraph{Edges.} $A = \mathrm{kNN}_{k=8}\!\left(|\mathbf{C}|\right)$, excluding the
diagonal: each ROI is connected to the eight regions with which its absolute correlation is
largest. Edge attributes are \textbf{binary}.

\[
  a_{ij} = \begin{cases}
    1 & \text{if } j \in \operatorname{top-8}_{j' \ne i} \left| c_{ij'} \right| \\
    0 & \text{otherwise}
  \end{cases}
\]

Two properties of this construction deserve to be stated explicitly rather than buried,
because Section~\ref{sec:gap-topology} identified them as the crux of a live dispute in the
field:

\begin{enumerate}
  \item The graph is a \emph{function of} the node features. It introduces no information
    that $X$ does not already contain.
  \item The construction is \emph{lossy in both directions}. Sparsification discards
    $192/199$ of each node's connections, and binarisation discards the magnitude of the
    eight that survive. A correlation of $0.81$ and one of $0.32$ become the same edge.
\end{enumerate}

\grey{The standard FC-graph construction gives the message-passing layer strictly less
information than the node features it operates on.}{%
If the encoder's value lies in message passing over $A$, it is passing messages over a
sparsified, binarised summary of $X$ --- while $X$ is available in full to any model that
simply pools it. This is the structural reason a no-encoder baseline is not a strawman but
the correct control, and it predicts the ablation result reported in
Chapter~\ref{ch:results} rather than rationalising it afterwards.}

%==============================================================================
\section{Stage 1: Static Graph Pretraining}
\label{sec:stage1}

\subsection{Graph Attention Autoencoder (GAAE)}

The encoder is a three-layer GATv2 network with FiLM conditioning. Each layer computes
attention coefficients over the neighbourhood,
\[
  \alpha_{ij} = \frac{\exp\!\left(\mathbf{a}^\top \mathrm{LeakyReLU}\!\left(\mathbf{W}[h_i \,\|\, h_j]\right)\right)}
                     {\sum_{k \in \mathcal{N}(i)} \exp\!\left(\mathbf{a}^\top \mathrm{LeakyReLU}\!\left(\mathbf{W}[h_i \,\|\, h_k]\right)\right)},
  \qquad
  h_i' = \sigma\!\left(\sum_{j \in \mathcal{N}(i)} \alpha_{ij} \mathbf{W} h_j\right),
\]
where GATv2's ordering of the non-linearity relative to the attention projection --- as
opposed to the original GAT --- is what makes the attention genuinely dynamic rather than
rank-one in the query.

Training minimises a two-term reconstruction loss: a topological term recovering the
adjacency and a feature term recovering the node features,
\[
  \mathcal{L}_{\mathrm{GAAE}} = \mathcal{L}_{\mathrm{topo}} + \lambda\,\mathcal{L}_{\mathrm{feat}},
\]
with cosine similarity as the reconstruction criterion. The visit embedding is the mean over
node latents, $z^{(t)} = \frac{1}{200}\sum_i z_i^{(t)}$, giving $d = 64$.

\subsection{Variational Graph Autoencoder (VGAE)}

The variational alternative encodes to a distribution rather than a point, optimising the
evidence lower bound
\[
  \mathcal{L}_{\mathrm{VGAE}} = \mathbb{E}_{q(Z|X,A)}\!\left[\log p(A \mid Z)\right]
    - \beta \, \mathrm{KL}\!\left(q(Z \mid X, A) \,\|\, p(Z)\right),
\]
with $p(Z) = \mathcal{N}(0, I)$. Both GCN and GAT message-passing variants were trained.
Posterior collapse --- the KL term driving $q$ to the prior so the latent carries nothing ---
is a live failure mode at this sample size, and the anti-collapse configurations in the
experiment registry exist to address it.

\subsection{What the pretext objective optimises for}

The objective above rewards a latent that can rebuild edges. It contains no term that
rewards separating converters from stable subjects, and the two are not the same
requirement (Section~\ref{sec:gap-pretext}). Between-subject variance that is
diagnostically informative but reconstruction-irrelevant has no gradient pushing it through
the 64-dimensional bottleneck. This is a stated design risk of the approach, not a
retrospective explanation.

%==============================================================================
\section{Stage 2: Longitudinal Trajectory Modelling (GE-LSTM / GE-GRU)}
\label{sec:gelstm-method}

\subsection{The recurrent core}

The trajectory classifier consumes the sequence $\{z^{(1)}, \ldots, z^{(T)}\}$ of per-visit
embeddings. At each step the embedding is concatenated with the elapsed-time feature before
entering the recurrent cell:
\[
  \tilde{z}^{(t)} = \left[\, z^{(t)} \;\|\; \Delta t^{(t)} \,\right] \in \mathbb{R}^{d+1},
  \qquad
  \Delta t^{(t)} = \frac{m^{(t)} - m^{(t-1)}}{m_{\max}},\quad \Delta t^{(1)} = 0,
\]
with $m^{(t)}$ the elapsed months at visit $t$ and $m_{\max}$ a fixed normaliser covering
the longest observed follow-up. Standard LSTM (or GRU) gate updates then run over
$\tilde{z}$, and the final hidden state $h^{(T)}$ is passed to a small MLP head producing a
conversion logit.

The recurrent input width is \emph{derived}, never hardcoded:
$\texttt{lstm\_input\_dim} = \texttt{embed\_dim} + 1$ when time conditioning is enabled.
This matters for the ablation of Section~\ref{sec:encoder-arms}, where removing the encoder
changes \texttt{embed\_dim} from 64 to 200 and the recurrent core must resize itself
accordingly.

\paragraph{On the strength of the time-conditioning claim.} Concatenating $\Delta t$ to the
input is the simplest form of time-awareness --- weaker than T-LSTM's memory decomposition
with elapsed-time decay, which modifies the cell state itself. This thesis does not claim
the mechanism is novel. Section~\ref{sec:delta-t-finding} additionally shows it cannot be
shown to contribute on DELCODE at all, because DELCODE's follow-up intervals are 90\%
identical; the mechanism is testable only on the external cohorts.

\subsection{Regularisation}

LayerNorm and dropout are applied in the recurrent and classifier stacks. When the encoder
is frozen, its BatchNorm layers are held in evaluation mode so that running statistics do
not drift with the downstream batches --- a detail that is easy to get wrong and that
silently changes a "frozen" encoder into a partially adaptive one.

%==============================================================================
\section{The Four Encoder Arms}
\label{sec:encoder-arms}

The central ablation of this thesis is controlled by a single configuration field,
\texttt{encoder\_init}, which selects among four arms. They form a chain of differences, so
that each adjacent pair isolates exactly one factor:

\begin{align*}
  \texttt{none}
    \;&\xrightarrow{\;+\ \text{graph encoder}\;}\; \texttt{random} \\[0.3em]
  \texttt{random}
    \;&\xrightarrow{\;+\ \text{reconstruction pretraining}\;}\; \texttt{pretrained\_frozen} \\[0.3em]
  \texttt{pretrained\_frozen}
    \;&\xrightarrow{\;+\ \text{task adaptation}\;}\; \texttt{pretrained\_finetuned}
\end{align*}

\begin{table}[htbp]
  \centering
  \small
  \begin{tabular}{@{}lccc p{0.40\linewidth}@{}}
    \toprule
    \textbf{Arm} & \textbf{Enc.} & \textbf{Wts.} & \textbf{Tr.} & \textbf{Isolates} \\
    \midrule
    \texttt{none}                 & no  & --- & --- & Whether a graph encoder earns its place at all \\
    \texttt{random}               & yes & no  & yes & Value of the pretraining, architecture held constant \\
    \texttt{pretrained\_frozen}   & yes & yes & no  & Reference --- the standard pipeline \\
    \texttt{pretrained\_finetuned}& yes & yes & yes & Value of adapting pretrained features to the task \\
    \bottomrule
  \end{tabular}
  \caption{The four encoder arms. Columns: whether an encoder is built (Enc.), whether
    pretrained GAAE weights are loaded (Wts.), and whether the encoder is trained (Tr.).
    All four arms share one seed, one fold assignment, one test manifest and one
    hyperparameter file, so the comparison is subject-for-subject.}
  \label{tab:encoder-arms}
\end{table}

\paragraph{The \texttt{none} arm.} No encoder module is constructed. The raw node features
are mean-pooled directly, so the visit embedding is the column-mean of the FC matrix --- a
200-dimensional global connectivity profile --- and the recurrent core widens to accept it.
Trainable parameters drop from $520{,}905$ to $31{,}169$.

\paragraph{A gradient-plumbing subtlety that makes the ablation meaningful.} The
visit-encoding routine historically wrapped every encoder call in \texttt{torch.no\_grad()}
and forced evaluation mode, which made the encoder a pure feature extractor \emph{regardless
of what \texttt{requires\_grad} said}. Under that behaviour the pre-existing
"unfrozen" configuration could not actually fine-tune anything, and the \texttt{random} arm
would have measured a frozen random encoder rather than a trained one. An explicit
\texttt{encoder\_grad} flag, defaulting to \texttt{False} and enabled only for the training
pass of the two arms that train their encoder, opens that context. Evaluation always encodes
under no-grad in eval mode. A regression test pins this behaviour, because without it the
ablation silently measures the wrong thing.

\edge{The ablation's interpretation table was pre-registered.}{%
Every possible outcome pattern --- including the ones unfavourable to the pipeline --- was
written down with its intended conclusion \emph{before} any of the four arms had been run.
The reader can therefore verify that the conclusion drawn in Chapter~\ref{ch:results} was
not selected after seeing the numbers. This is standard in clinical trials and vanishingly
rare in methods papers.}

%==============================================================================
\section{Baseline Classifiers \& Static Formulations}
\label{sec:static-baselines}

Three non-recurrent comparators are trained on identical splits.

\begin{description}[leftmargin=2.6cm,style=nextline]
  \item[GEC] Graph Encoder Classifier --- precomputes frozen per-visit embeddings and trains
    a deep MLP head (hidden widths $[256, 128, 64]$ with BatchNorm) over the visit set,
    with an explicit visit-count mask as its only length signal.
  \item[GEP] Graph Embedding Perceptron --- a shallower pooled-embedding MLP over the same
    frozen embeddings.
  \item[Raw FC] A classifier applied directly to the connectivity features with no graph
    encoder, serving as the representation floor.
\end{description}

\risk{GEC's explicit visit-count mask is the one place a length shortcut could enter.}{%
The recurrent models never receive the visit count as a feature; GEC's
\texttt{append\_visit\_mask} does supply length information explicitly. Given that
converters in this cohort are followed longer (Section~\ref{sec:gap-attrition}), any GEC
result must be read against the within-subject slope diagnostics before it is compared with
the recurrent arms on equal terms.}

%==============================================================================
\section{Class Imbalance, Thresholds \& Evaluation Protocol}
\label{sec:thresholds}

\paragraph{Imbalance.} Converters are the minority class in every split. Training uses
\texttt{BCEWithLogitsLoss} with $\texttt{pos\_weight} = n_{\mathrm{neg}}/n_{\mathrm{pos}}$,
applied identically across all arms of any comparison --- changing it in one arm only would
confound the comparison it is meant to support.

\paragraph{Threshold selection --- the project's hard rule.} A binary decision requires a
threshold, and where that threshold comes from is a leakage question, not a tuning question.
Two validation-derived rules are supported: the $F_1$-optimal point, and Youden's
$J = \mathrm{sensitivity} + \mathrm{specificity} - 1$. Both are computed on out-of-fold
\emph{validation} predictions and then applied unchanged to the test set.

This is enforced rather than documented. The evaluation entry point raises when no explicit
threshold is supplied:

\begin{quote}\small\ttfamily
  evaluate\_classifier requires an explicit threshold (typically the validation-derived
  best\_threshold). Choosing a threshold from test metrics would leak.
\end{quote}

A silent default of $0.5$ would be the single easiest way to leak test information into a
reported sensitivity/specificity pair, and the loud failure is the intended behaviour.

\paragraph{Reported metrics.} AUC-ROC is the primary metric, being threshold-free.
Sensitivity, specificity, $F_1$ and balanced accuracy are reported at the
validation-derived operating point. Cross-validation results are reported as a distribution
over folds; held-out test results as a point estimate with the cohort size stated alongside,
since at $N_{\text{test}} = 34$ the estimate's own uncertainty is large.

\paragraph{Reference floors.} Every results table is anchored against the metadata-only
baseline --- age, sex and visit timing through a gradient-boosted classifier, which attains
CV AUC $0.6157 \pm 0.0653$ and test AUC $0.4929$ --- and against the raw-FC representation
floor. An arm that fails to clear the metadata floor is not evidence about encoders.

%==============================================================================
\section{Scope of the Time-to-Event Extension (PROGNOSER)}
\label{sec:prognoser-scope}

Section~\ref{sec:survival} argued that binary conversion prediction is a compromise: the
label is a censored observation treated as a class, and a clinician needs a time, not a
flag. A survival extension consuming the same graph latents was therefore planned, modelling
the hazard directly and reporting the concordance index and time-dependent AUC.

\risk{The survival extension is scoped, not delivered, and this thesis says so rather than
implying otherwise.}{%
What exists is a Kaplan--Meier baseline and a Cox model on \emph{clinical} covariates only.
No imaging-based survival head has been trained, and no C-index for the graph-latent
pipeline exists. Presenting the survival component as a result would misrepresent the state
of the work; it is stated here as motivation and carried to Chapter~\ref{ch:conclusion} as
future work. The classification results stand on their own without it.}
